\subsection{Characterizations of Sobolev Spaces}\label{sec:sobolev-characterizations}

The basic definition of a Sobolev space describes it as a subspace of $L^p(\Omega)$. Of course, there is much more to be said, and in this section we describe some of the most important ways that functions in a Sobolev space can be characterized.

In most of this section, we shall confine our discussions to the case $p=2$. Many of the results we discuss have analogues for general $p$, for which we refer to the literature.

\subsubsection{Some Comments on the Domain $\Omega$}\label{subsec:sobolev-domain-boundary-hypotheses}

The answers to a number of questions about Sobolev spaces depend on assumptions on the regularity of the boundary of $\Omega$.\footnote{In this rather short treatment of Sobolev spaces, we have chosen to avoid most questions of boundary smoothness. For a more complete study of the subject we recommend the paper of Fraenkel \cite{Fr}.} Most of the time, we shall assume a smooth boundary. Specifically, we make the following definition.

\begin{Definition}
\label{def:c-k-domain-boundary}
We say that $\Omega$ is of class $C^k$, $k\geq1$, if every point on $\partial\Omega$ has a neighborhood $N$ so that $\partial\Omega\cap N$ is a $C^k$-surface and, moreover, $\Omega\cap N$ is ``on one side'' of $\partial\Omega\cap N$.
\end{Definition}

If $\Omega$ is a bounded domain, i.e., connected, the last assumption is redundant; cf. Remark~\ref{rem:boundary-one-side-jordan}.

There are two classes of problems in applications, where nonsmooth domains are relevant:\begin{enumerate}
\item domains with corners, and
\item free boundary problems where $\Omega$ is a priori unknown.
\end{enumerate}It turns out that in fact many results on Sobolev spaces do not require a smooth boundary. Instead, various geometric conditions such as the ``segment property'' and ``cone property'' (cf. \cite{Fr}) need to be assumed. We shall not discuss these conditions here, but we shall state some results for Lipschitz domains.

\begin{Definition}
\label{def:lipschitz-domain}
We say that $\Omega$ is \textbf{Lipschitz} if every point on $\partial\Omega$ has a neighborhood $N$ such that, after an affine change of coordinates (translation and rotation), $\partial\Omega\cap N$ is described by the equation $x_m=\phi(x_1,\ldots,x_{m-1})$, where $\phi$ is \textbf{uniformly Lipschitz continuous}. Moreover, $\Omega\cap N$ is on one side of $\partial\Omega\cap N$, e.g., $\Omega\cap N=\{\bm{x}\in N\mid x_m<\phi(x_1,\ldots,x_{m-1})\}$.
\end{Definition}

If $\Omega$ is unbounded, then, in addition to smoothness conditions on $\partial\Omega$, one needs to impose conditions which say that $\Omega$ is well behaved at infinity. We shall not give a general discussion of such conditions, and many results will be stated only for the case when $\partial\Omega$ is bounded.

Finally, we define a characterization of the domain $\Omega$ that will be very useful as a concise technical hypothesis.

\begin{Definition}
\label{def:sobolev-k-extension-property}
We say that $\Omega$ has the \textbf{$k$-extension property} if there is a bounded linear mapping $E:H^k(\Omega)\to H^k(\mathbb{R}^m)$ such that $Eu|_{\Omega}=u$ for every $u\in H^k(\Omega)$.
\end{Definition}

It is of course trivial that, conversely, the restriction of every function in $H^k(\mathbb{R}^m)$ is in $H^k(\Omega)$. The extension property will be investigated in a later subsection; it turns out that bounded Lipschitz domains have the extension property for every $k$.
